\section{Related Work}
\label{sec:related}

A substantial body of work investigates formal verification for Rust. One line of research focuses on automated verification.
Prusti~\cite{astrauskas2019leveraging} and Creusot~\cite{denis2022creusot} translate Rust programs and specifications into verification conditions discharged by SMT-based backends.
Verus~\cite{lattuada2023verus} integrates specifications and ghost code with executable Rust, while Kani~\cite{kani} applies bounded model checking~\cite{vanhattum2022verifying}.
These tools can discharge many proof obligations with limited user interaction.
However, SMT-based verification is restricted to properties expressible within its supported specification and solver fragments, while bounded model checking explores executions only up to a chosen bound.
This limits the flexibility and extensibility of proofs that require custom abstractions or reasoning outside these fragments.
Recent work uses AI agents to reduce the effort of interactive formal development, but still assumes that a formal program model is already available in the target prover~\cite{fang2026trustworthy}.

Another line of work constructs Rust program models for interactive proof assistants.
Electrolysis~\cite{ullrich2016simple} introduced a functional translation of a Rust subset into Lean, and hax~\cite{bhargavan2024hax} generates models for proof-oriented languages including F* and Coq.
Aeneas~\cite{aeneas_project,ho2022aeneas} further develops this direction by translating a substantial safe-Rust fragment into pure functions and providing backends for F*, Coq, HOL4, and Lean.
RustBelt~\cite{jung2017rustbelt} and RefinedRust~\cite{gaher2024refinedrust} instead use separation logic in Coq to reason about Rust memory and ownership explicitly.
In contrast to these existing target choices, \texttt{Rust2Isabelle} translates Aeneas's functional representation into Isabelle/HOL theories, enabling Rust programs to be studied using Isabelle's higher-order logic and interactive proof facilities.

Isabelle/HOL is an attractive verification target because it combines an expressive logic with mature libraries and proof automation.
AutoCorres~\cite{greenaway2014don} translates C programs into abstract monadic models for Isabelle proofs.
AutoCorrode~\cite{autocorrode_aws} provides separation-logic infrastructure for \(\mu\)Rust, a restricted Rust-like language embedded in Isabelle/HOL, and therefore requires programs to be modeled in that language.
Neither approach automatically translates Rust source code into Aeneas's value-based functional representation.
\texttt{Rust2Isabelle} fills this missing connection by extending the Aeneas toolchain with an Isabelle/HOL backend and the semantic support required by the generated theories.